% B. Análisis de necesidades y estudio del estado del arte (Guía apdo. 3.B)
% Extensión orientativa: 5-7 páginas. Capítulo con mayor densidad de citas IEEE.
\chapter{Análisis de necesidades y estado del arte}
\label{cap:estado-arte}

\section{Análisis de necesidades}

El sistema debe satisfacer los siguientes requisitos, derivados del escenario de aplicación:

\begin{itemize}
    \item \textbf{Captura multivariable.} Un único tipo de señal no basta para discriminar
        modos de fallo: la vibración revela desequilibrios mecánicos, la temperatura revela
        pérdidas de rendimiento térmico y la firma acústica revela anomalías del régimen de
        giro.
    \item \textbf{Resiliencia mecánica.} El nodo opera acoplado a un elemento en vibración
        permanente, lo que degrada conexiones y cableado. El sistema debe seguir midiendo
        ante fallos físicos transitorios.
    \item \textbf{Operación autónoma prolongada.} La construcción del modelo de
        comportamiento nominal exige campañas de captura de días, sin supervisión.
    \item \textbf{Independencia de la nube.} El diagnóstico debe poder resolverse en el
        nodo, con la infraestructura externa limitada al entrenamiento y al análisis.
    \item \textbf{Coste contenido.} La solución debe construirse con componentes de bajo
        coste para resultar aplicable a activos de valor moderado.
\end{itemize}

\section{Mantenimiento predictivo basado en señales físicas}

La monitorización de condición de máquina rotativa se apoya en un conjunto de indicadores
consolidado, cuya elección responde a que cada uno resulta sensible a una familia distinta de
mecanismos de fallo~\cite{randall2011}. Esa correspondencia entre indicador y mecanismo es la
que justifica la selección de características descrita en el capítulo~\ref{cap:diseno}.

El valor eficaz de la componente alterna de la aceleración constituye el indicador de
severidad global. Su evolución es monótona con la degradación, lo que lo hace adecuado para
seguimiento de tendencia, y existen criterios normativos que relacionan su magnitud con
estados admisibles de la máquina según su clase~\cite{iso20816}. Su limitación es la
contrapartida de esa robustez: al promediar sobre toda la ventana, atenúa los eventos poco
frecuentes.

La kurtosis cubre precisamente ese hueco. Al ser el cuarto momento normalizado de la
distribución de amplitudes, cuantifica la impulsividad de la señal y no su magnitud, con lo
que responde a la aparición de impactos aislados antes de que estos alteren el valor
eficaz~\cite{randall2011}. Es el indicador clásico del deterioro incipiente de rodamientos, y
su interpretación se apoya en dos valores de referencia analíticos: 3 para ruido gaussiano y
1,5 para una señal senoidal pura. El par formado por el valor eficaz y la kurtosis discrimina
así entre mecanismos que un único indicador confundiría: un desequilibrio de masa eleva el
primero manteniendo el segundo próximo al de una senoide, mientras que un rodamiento en
deterioro incipiente apenas altera el primero y sin embargo eleva el segundo.

El análisis espectral aporta la dimensión que los indicadores temporales no pueden ofrecer: la
localización en frecuencia. La componente correspondiente a la frecuencia de giro y la
relación entre esta y sus armónicos permiten distinguir desequilibrio de masa, desalineación
y aflojamiento del anclaje, mecanismos que producen niveles globales semejantes pero
distribuciones espectrales distintas~\cite{randall2011}.

Procede señalar una condición que la literatura recoge y que resultó determinante en este
trabajo: la banda de frecuencias efectivamente utilizable no la fija el sensor sino su
montaje. La fijación del acelerómetro sobre la superficie a medir constituye un sistema
elástico cuya resonancia limita el intervalo en que la medida es fiel, y existen criterios
normalizados para el montaje precisamente por esa razón~\cite{iso5348}. El
capítulo~\ref{cap:resultados} documenta la magnitud de ese efecto en el montaje empleado.

La monitorización acústica se plantea como modalidad complementaria y no como alternativa. Su
interés reside en que un micrófono capta la señal radiada al aire sin requerir contacto
mecánico, lo que la hace sensible a fenómenos que la vibración estructural no transmite bien,
como la obstrucción del caudal de aire. Su inconveniente, contrastado experimentalmente en
este trabajo, es que un micrófono solidario a la estructura recibe la vibración por conducción
y deja de medir el aire.

\section{Inteligencia en el borde de la red}

La ejecución de modelos de aprendizaje automático sobre microcontroladores de bajo consumo
constituye un campo con denominación propia, habitualmente designado como \textit{TinyML}, y
su desarrollo responde a restricciones que no aparecen en los entornos de cómputo
convencionales: memoria del orden de decenas o centenares de kilobytes, ausencia de sistema
operativo de propósito general y presupuesto energético acotado~\cite{warden2019}. Las
herramientas que han hecho practicable ese despliegue, en particular los entornos de
ejecución de modelos reducidos para microcontroladores, prescinden de asignación dinámica de
memoria y de gran parte del repertorio de operadores de sus equivalentes de
servidor~\cite{david2021}.

Conviene precisar el alcance de esas herramientas, porque condiciona la solución que se adopta
en el capítulo~\ref{cap:resultados}. Su repertorio de operadores corresponde a
\textbf{redes neuronales}: convoluciones, capas densas y funciones de activación. Un modelo que
no pertenezca a esa familia —una envolvente estadística, un método basado en vecindad, un
conjunto de árboles— no es ejecutable en ellas, de modo que su despliegue sobre el
microcontrolador exige programarlo directamente. La disponibilidad de un entorno de inferencia
no es por tanto un requisito del cómputo en el borde, sino una facilidad aplicable a una
familia concreta de modelos.

El contraste con una arquitectura orientada a la nube resulta cuantificable en el caso que
nos ocupa, y no es una cuestión de preferencia. Transmitir la señal de vibración en bruto de
los tres ejes a la frecuencia que exige el análisis frecuencial supone del orden de 3000
valores por segundo de forma continua. Extraer las características en el nodo y transmitir
únicamente estas reduce el volumen a un mensaje cada \mbox{30 s}, tres órdenes de magnitud por
debajo. La diferencia no afecta solo al ancho de banda: condiciona la autonomía del nodo, el
coste de almacenamiento y, sobre todo, la dependencia operativa de un enlace externo.

Esa última consideración es la que motiva el planteamiento de este trabajo. Un sistema de
diagnóstico cuya lógica reside en un servidor remoto deja de diagnosticar cuando el enlace
falla, circunstancia que en un entorno industrial no es excepcional. Situar la inferencia en
el nodo convierte la infraestructura externa en soporte de entrenamiento y análisis en lugar
de requisito de operación.

Conviene precisar que el planteamiento no exige necesariamente una red neuronal. Cuando el
problema admite una formulación estadística, el modelo embarcado puede reducirse a un
conjunto de parámetros de unos pocos kilobytes, sin la infraestructura de un entorno de
ejecución de aprendizaje automático. El capítulo~\ref{cap:diseno} desarrolla esta alternativa
y su justificación.

\section{Detección de anomalías no supervisada}

La formulación del problema de detección responde a una asimetría en la disponibilidad de
datos que es característica del mantenimiento predictivo. De un activo en servicio se obtiene
con facilidad un registro extenso de su comportamiento normal, mientras que los ejemplos de
fallo son escasos, difíciles de etiquetar y, cuando se obtienen de forma deliberada, costosos.
Un clasificador supervisado requeriría un número de ejemplos de cada modo de fallo que no es
alcanzable en este contexto, de modo que el planteamiento adecuado es el de detección de
anomalías con una sola clase: el modelo se ajusta exclusivamente sobre el estado nominal y
señala como anómala toda observación suficientemente alejada de la región que ha aprendido a
describir.

Entre los métodos aplicables destaca por su idoneidad el basado en árboles de aislamiento, que
en lugar de modelar la densidad de la distribución normal explota que las observaciones
anómalas resultan más fáciles de separar del resto mediante particiones
aleatorias~\cite{liu2008}. Su interés práctico reside en un coste computacional reducido y en
que no requiere estimar una densidad en un espacio de características de dimensión elevada,
condición que otros métodos encuentran problemática.

La alternativa clásica formula el problema como la estimación del soporte de una distribución
de dimensión elevada, buscando la frontera que encierra la fracción especificada de los datos
de entrenamiento~\cite{scholkopf2001}. Su comportamiento depende de la elección del núcleo y
de sus parámetros, y su coste crece con el número de observaciones, lo que la hace menos
conveniente cuando el conjunto de referencia es extenso.

Procede añadir dos métodos que el desarrollo de este trabajo mostró pertinentes y que la
formulación inicial no contemplaba. El primero estima la densidad local de cada observación y la
compara con la de su vecindad, con lo que declara atípico aquello que reside en una región
menos poblada que su entorno inmediato~\cite{breunig2000}. Su interés en este caso reside en
que no presupone una forma concreta para la distribución del estado nominal, circunstancia
favorable cuando ese estado agrupa condiciones de operación heterogéneas, como son los ciclos
de marcha de duración desigual de un compresor sometido a uso doméstico.

El segundo ajusta una envolvente elipsoidal mediante una estimación robusta de la matriz de
covarianza, obtenida a partir del subconjunto de observaciones de menor determinante, lo que
impide que unas pocas observaciones atípicas del propio conjunto de ajuste desplacen la
frontera~\cite{rousseeuw1999}. Su ventaja para el objetivo de este trabajo es que la regla de
decisión resultante es una forma cuadrática con parámetros precalculados, ejecutable en el nodo
con una ocupación de memoria del orden de centenares de bytes.

Una tercera vía la constituyen los autocodificadores, que se entrenan para reconstruir la
señal de entrada y emplean el error de reconstrucción como puntuación de anomalía. Su ventaja
es que encajan de forma natural en los entornos de ejecución para microcontroladores, dado
que son redes neuronales; su inconveniente, un coste de entrenamiento y una necesidad de datos
superiores a los de los dos métodos anteriores.

En los tres casos, la decisión final descansa sobre un umbral aplicado a una puntuación
escalar. Fijar ese umbral sobre un percentil de la distribución de puntuaciones del propio
conjunto nominal hace explícita y ajustable la tasa de falsos positivos admitida, lo que
resulta preferible a un valor arbitrario.

\section{Síntesis y aportación de este trabajo}

Los trabajos revisados abordan por separado las tres cuestiones que este proyecto integra: los
indicadores de condición sobre señal de vibración, el despliegue de inferencia en
microcontroladores y la detección de anomalías con una sola clase. La combinación de las tres
sobre una plataforma de coste reducido está menos explorada, y es donde este trabajo pretende
aportar.

Existe además una circunstancia que la literatura consultada tiende a no abordar, y que en la
práctica de este proyecto resultó ser la determinante. Los métodos de monitorización de
condición se formulan asumiendo instrumentación de calidad y adquisición fiable: acelerómetros
acoplados de forma rígida y una cadena de medida que entrega las muestras que se le piden. Al
trasladar esos métodos a hardware de prototipado adherido a una máquina en vibración, ninguna
de las dos premisas se cumple. El bus de comunicación con el sensor falla de forma
intermitente, la fijación del acelerómetro introduce una resonancia que limita la banda
utilizable, y ambos efectos degradan precisamente los indicadores más sensibles.

La aportación de este trabajo se sitúa por tanto en dos planos. En el primero, la integración
de tres modalidades de sensado de bajo coste con extracción de características e inferencia en
el propio nodo. En el segundo, y con mayor valor práctico, la documentación cuantificada de
las limitaciones que impone ese hardware y de los mecanismos que permiten operar con ellas: la
validación de la integridad del dato en la adquisición, la instrumentación que hace medible la
calidad de cada medida, y la acotación explícita de la banda de frecuencias en que las
características resultan interpretables. El capítulo~\ref{cap:conclusiones} valora en qué
medida esas limitaciones condicionan el alcance alcanzado.
